\documentclass[11pt]{article}

\usepackage[french, english]{babel}

\usepackage[T1]{fontenc}

\usepackage{amsfonts}
\usepackage{amsmath}
\usepackage{amssymb}

\usepackage{graphicx}
\usepackage{multirow}
\usepackage{multicol}
\usepackage{xcolor}
\usepackage[shortlabels]{enumitem}

\usepackage[margin=2cm]{geometry}

\usepackage{mathtools}
\usepackage{siunitx}
\usepackage{braket}

\begin{document}
\begin{multicols}{2}

\section{\underline{Description statistique}}
\textbf{Micro-état :} ex: position et vitesse de chaque molécule d'un gaz.\\
\textbf{Macro-état :} ex: pression, volume, température d'un gaz.\\
\textbf{Postulat fondamental :} Tous les micro-états compatibles avec un macro-état sont équiprobables :
\begin{equation*}
  P_i = \frac{1}{\Omega (E, V, N)}
\end{equation*}\\
où $\Omega (E, V, N)$ est le nombre de micro-états accessibles.\\
\textbf{Entropie :}
\begin{equation*}
  \boxed{S(E, V, N) = k_B \ln \Omega (E, V, N)}
\end{equation*}\\
(Trucs : Système isolé : $S$ max.\\
Travailler sur $\ln\Omega$)\\

\section{\underline{Ens. micro. (E, V, N) constants}}
\textbf{Distribution :} Postulat fondamental.\\
\textbf{Température :}
\begin{equation*}
  \frac{1}{T} = \left( \frac{\partial S}{\partial E} \right)_{V, N}
\end{equation*}\\
\textbf{Pression :}
\begin{equation*}
  \frac{P}{T} = \left( \frac{\partial S}{\partial V} \right)_{E, N}
\end{equation*}\\
\textbf{Potentiel chimique :}
\begin{equation*}
  \frac{\mu}{T} = -\left( \frac{\partial S}{\partial N} \right)_{E, V}
\end{equation*}\\
\textbf{Exemples :} Gaz parfait et système de spins.\\
(Truc : Bien indentiifier les constantes avant tout.)\\

\section{\underline{Ens. can. (T, V, N) constants}}
\textbf{Distribution :}
\begin{equation*}
  \boxed{P_i = \frac{1}{Z} e^{-\beta E_i}}, \boxed{Z = \sum_i e^{-\beta E_i}} \quad \text{avec} \quad \beta = \frac{1}{k_B T}
\end{equation*}
\textbf{Grandeurs thermodynamiques :}
\begin{equation*}
  \braket{E}=-\dfrac{\partial}{\partial\beta} \ln Z \quad S = k_B(\ln Z + \beta \braket{E})\quad C_V = \left( \frac{\partial \braket{E}}{\partial T} \right)_{V, N}
\end{equation*}
\textbf{Relations thermodynamiques :}
\begin{equation*}
  S = -\left( \frac{\partial F}{\partial T} \right)_{V, N}  P = -\left( \frac{\partial F}{\partial V} \right)_{T, N}  \mu = \left( \frac{\partial F}{\partial N} \right)_{T, V}
\end{equation*}\\
\textbf{Exemples :} Oscillateurs harmoniques, spins, gaz parfait.\\
(Truc : Si t'as Z, tu peux avoir F, S, E, $C_V$.)\\

\section{\underline{Ens. gran. (T, V, $\mu$) constants}}
\textbf{Distribution :}
\begin{equation*}
  \boxed{P_i = \frac{1}{\mathcal{Z}} e^{-\beta (E_i - \mu N_i)}} \quad \text{avec} \quad \boxed{\mathcal{Z} = \sum_i e^{-\beta (E_i - \mu N_i)}}
\end{equation*}
\textbf{Grandeurs thermodynamiques :}
\begin{equation*}
  \braket{E} = -\frac{\partial}{\partial \beta} \ln \mathcal{Z} \quad \braket{N} = \frac{1}{\beta} \left( \frac{\partial}{\partial \mu} \ln \mathcal{Z} \right)_{T, V}
\end{equation*}
\begin{equation*}
  S = k_B (\ln \mathcal{Z} + \beta \braket{E} - \beta \mu \braket{N})
\end{equation*}
\textbf{Relations thermodynamiques :}
\begin{equation*}
  S = -\left( \frac{\partial \Omega}{\partial T} \right)_{V, \mu}  P = -\left( \frac{\partial \Omega}{\partial V} \right)_{T, \mu} \braket{N} = -\left( \frac{\partial \Omega}{\partial \mu} \right)_{T, V}
\end{equation*}\\
\textbf{Exemples :} Gaz parfait classique, photons ($\mu=0$), bosons/fermions.\\
(Truc : Si N fluctue $\rightarrow$ Grand canonique obligé)\\

\section{\underline{Autres}}
\textbf{Gaz parfait :}
\begin{equation*}
  PV = N k_B T \quad \braket{E}=\frac{3}{2}Nk_BT  
\end{equation*}
\textbf{Enthalpie :}
\begin{align*}
  H &= E + PV\\
  &dH = dE + P dV + V dP\\
  &= T dS + V dP\\
  &\left( \frac{\partial T}{\partial P} \right)_S = -\left( \frac{\partial V}{\partial S} \right)_P
\end{align*}
\textbf{Énergie libre de Helmholtz :}
\begin{align*}
  F &= E - TS\\
  dF &= dE - T dS - S dT\\
  &= -S dT - P dV\\
  &\left( \frac{\partial S}{\partial V} \right)_T=\left( \frac{\partial P}{\partial T} \right)_V
\end{align*}
\textbf{Énergie libre de Gibbs :}
\begin{align*}
  G &= H - TS\\
  &=E+PV-TS\\
  &=F+PV-TS\\
  dG &= -SdT+VdP\\
  &\left( \frac{\partial S}{\partial P} \right)_T = -\left( \frac{\partial V}{\partial T} \right)_P
\end{align*}
\textbf{Grand potentiel :}
\begin{equation*}
  \Omega = \braket{E}-TS-\mu N
\end{equation*}
\textbf{Capacité thermique :}
\begin{equation*}
  C_V = \left( \frac{\partial \braket{E}}{\partial T} \right)_{V, N} \quad C_P = \left( \frac{\partial \braket{H}}{\partial T} \right)_{P, N}
\end{equation*}
\textbf{Lois de la thermodynamique :}
\begin{itemize}
  \item $\boxed{dU = PdV - TdS}$ (1ère loi)\\
        Cas utiles : \begin{itemize}
          \item Isotherme : $dU=0 \Rightarrow \delta Q = \delta W$
          \item Adaiabatique : $\delta Q = 0 \Rightarrow dU = \delta W$
          \item Isochore : $\delta W = 0 \Rightarrow dU = \delta Q$
        \end{itemize}
        (Truc : Préciser le signe du travail)
  \item $\boxed{\Delta S_{univers} \geq 0}$\\
        $dS = \frac{\delta Q_{rev}}{T}$\\
        (Truc : Calculer $\Delta S$ le long d'un processus réversible)
  \item $\boxed{S \xrightarrow[T\rightarrow 0]{} S_0 (cste)}$\\
        (Truc : $S_0$ est la valeur de l'entropie à zéro absolu)
\end{itemize}

\subsection{Machines thermiques}
\begin{equation*}
  \boxed{W = Q_h - Q_c}
\end{equation*}
\textbf{Rendement}
\begin{equation*}
  \boxed{\eta = \frac{W}{Q_h} = 1 - \frac{Q_c}{Q_h} \leq 1 - \frac{T_c}{T_h}}
\end{equation*}
(Truc : Convertir en kelvins pour le rendement)\\
\textbf{Coefficient de performance (réfrigérateur)}
\begin{equation*}
  \boxed{COP = \frac{Q_c}{W} = \frac{1}{\frac{Q_h}{Q_c} - 1} \leq \frac{1}{\frac{T_h}{T_c} - 1}}
\end{equation*}\\
\textbf{Coefficient de performance (Pompe à chaleur)}
\begin{equation*}
  \boxed{COP = \frac{Q_h}{W} = \frac{1}{1 - \frac{Q_c}{Q_h}} \leq \frac{1}{1 - \frac{T_c}{T_h}}}
\end{equation*}
\textbf{Entropie dans une machine thermique réversible}
\begin{equation*}
  \boxed{\Delta S = -\frac{Q_h}{T_h} + \frac{Q_c}{T_c} \geq 0}
\end{equation*}

\subsection{Aides}
\textbf{Pour les ensembles :}
\begin{itemize}
  \item \begin{itemize}
          \item Microcanonique : Identifier $E$, $V$, $N$ constants.
          \item Canonique : Identifier $T$, $V$, $N$ constants. Trouver $Z$.
          \item Grand canonique : Identifier $T$, $V$, $\mu$ constants. Trouver $\mathcal{Z}$.
        \end{itemize}
  \item Trouver Z
  \item Prendre le log
  \item Dériver pour trouver les grandeurs thermodynamiques
  \item Faire l'interprétation physique
\end{itemize}



\end{multicols}
\end{document}